\section{The \texttt{normalize\_extremes} Path: An Architectural Residue}
\label{app:evt_path}

This appendix documents an architectural pathway in the library
that the empirical results discourage: the Hill-estimator-driven
data-layer rescaling implemented by
\texttt{sert\_xlearner.core.orthogonalization.normalize\_extremes}.
We describe it here because the code remains in the public library
for reproducibility of earlier experiments; users must pass
\texttt{tail\_threshold} and \texttt{tail\_alpha} explicitly to
activate it, and we recommend they do not.

\paragraph{What the code does.}
\texttt{estimate\_tail\_parameters} implements a standard Hill
estimator on residuals: for a top-percentile threshold $t$ and
exceedances $\{|r_i| : |r_i| > t\}$, the estimator returns
$\hat\alpha = 1 / \overline{\ln(|r_i| / t)}$, the inverse of the
sample Pareto-tail slope. \texttt{normalize\_extremes} then
transforms Phase 2 pseudo-outcomes as
$D_i \mapsto D_i / t^{\hat\alpha}$ for $|D_i| > t$ and leaves bulk
pseudo-outcomes unchanged. The intent was a tail-scale normaliser
that would stabilise the MCMC gradient regardless of whether the
tail carried signal or contamination.

\paragraph{What empirical evaluation shows.}
The tail-heterogeneous probe of
Section~\ref{sec:experiments:tail_signal} activates this pathway on
a DGP where the tail is explicitly signal: whales have
$\tau_{\text{whale}} = 10$ while the bulk has $\tau_{\text{bulk}} = 2$.
The \texttt{normalize\_extremes} operator divides the whale
pseudo-outcomes by $t^\alpha$ --- which for Hill-estimated
$(t, \hat\alpha)$ on this DGP is approximately $(5, 1.5)$, a
division factor of $\approx 11$. The whale signal collapses to the
bulk scale, the posterior concentrates on a near-zero effect, and
subgroup coverage drops to zero. Across all six (basis,
likelihood) combinations of Table~\ref{tab:tail_signal}, activating
\texttt{normalize\_extremes} degrades both mixed and whale
$\eATE$ and kills subgroup coverage.

\paragraph{Why it fails.}
The operator is contamination-directed by construction: scaling
tail pseudo-outcomes \emph{down} corresponds to a Bayesian prior
belief that tail residuals are magnified noise rather than a
heterogeneous effect. When that belief is correct, the Welsch
likelihood already handles the case (and handles it better,
because Welsch's suppression is smooth rather than a hard
threshold). When the belief is wrong, as in the tails-as-signal
regime, the operator removes signal that would otherwise drive
the posterior. It is, in short, the wrong architectural layer:
tail modelling belongs in the likelihood or the prior, not as a
deterministic data transformation.

\paragraph{What a correct EVT-Bayesian likelihood would look like.}
A principled replacement would place a mixture likelihood at the
Phase 3 layer: a Welsch component for bulk residuals plus a
generalised Pareto tail component with shape parameter estimated
jointly with $\beta$. The Hill estimator would elicit an
informative prior on the Pareto shape rather than deterministically
rescale the data. Such a model is straightforward to implement in
NumPyro but introduces identifiability concerns that we have not
yet investigated at scale. We flag this as future work
(Section~\ref{sec:discussion}); the architectural residue in the
current library is neither silently removed nor silently
activated, but explicitly documented here as a known-bad
configuration.
